\section{压缩映射}
\label{sec:mmap-contraction}

许多方程可以写成$x=T(x)$的形式，满足它的点称为$T$的\term{不动点}（fixed point）。最朴素的解法是反复代入：任取初值$x^{(0)}$，计算$x^{(1)}=T(x^{(0)})$，$x^{(2)}=T(x^{(1)})$，……，希望迭代列趋于一个解。例如解$x=(x+2)/3$时，从$0$出发得到$2/3,\ 8/9,\ 26/27,\ldots$，看上去趋于$1$。本节说明这种做法什么时候一定成功：$T$把空间映回自身，按统一的比例缩小距离，并且空间完备。它把Lipschitz误差控制与完备性结合起来，不仅给出解的存在与唯一，还给出迭代误差的界。

\begin{definition}[压缩映射]
\label{def:mmap-contraction}
设$(X,d)$是度量空间。映射$T:X\to X$若存在常数$q\in[0,1)$，使对所有$x,y\in X$有
\begin{equation}
d(T(x),T(y))\leq q\,d(x,y),
\label{eq:mmap-contraction}
\end{equation}
则称$T$为\term{压缩映射}（contraction mapping），$q$称为\term{压缩常数}（contraction constant）。
\end{definition}

定义中有三个要求，后面的反例会逐一说明它们都不能去掉：$T$的值仍在$X$中，所以迭代可以无限进行；同一个$q<1$对所有点对有效，而不只是每对点各自缩小距离；$q<1$严格成立。压缩映射以$q$为Lipschitz常数，因而连续。

\begin{theorem}[压缩映射定理]
\label{thm:mmap-banach-fixed-point}
设$(X,d)$是非空完备度量空间，$T:X\to X$是压缩常数为$q$的压缩映射。则$T$有唯一的不动点$x^\ast$；对任意初值$x^{(0)}\in X$，迭代$x^{(k+1)}=T(x^{(k)})$收敛于$x^\ast$，并且对$k\geq1$，
\begin{equation}
d(x^{(k)},x^\ast)\leq\frac{q^k}{1-q}\,d(x^{(1)},x^{(0)}),\qquad
d(x^{(k)},x^\ast)\leq\frac q{1-q}\,d(x^{(k)},x^{(k-1)}).
\label{eq:mmap-contraction-error-bounds}
\end{equation}
\end{theorem}

下面三个小节依次证明迭代收敛于不动点、不动点唯一，以及两个误差界。

\subsection{迭代收敛}
\label{sec:mmap-iteration-convergence}

先估计相邻两次迭代之差。由$x^{(k+1)}=T(x^{(k)})$与$x^{(k)}=T(x^{(k-1)})$，压缩条件给出$d(x^{(k+1)},x^{(k)})\leq q\,d(x^{(k)},x^{(k-1)})$，反复使用得
\begin{equation*}
d(x^{(k+1)},x^{(k)})\leq q^k\,d(x^{(1)},x^{(0)}).
\end{equation*}
每一步的移动量按几何级数缩小。

相邻项接近还不够，柯西列要求任意两个后项接近。对$m>k$，沿途使用三角不等式，再对几何级数求和：
\begin{equation}
d(x^{(m)},x^{(k)})\leq\sum_{j=k}^{m-1}d(x^{(j+1)},x^{(j)})\leq\sum_{j=k}^{m-1}q^j\,d(x^{(1)},x^{(0)})\leq\frac{q^k}{1-q}\,d(x^{(1)},x^{(0)}).
\label{eq:mmap-contraction-cauchy}
\end{equation}
右端与$m$无关，并且$q<1$时随$k\to\infty$趋于零，所以$(x^{(k)})$是柯西列。

由$X$完备，迭代列收敛到某个$x^\ast\in X$。再用$T$的连续性：$T(x^\ast)=\lim_kT(x^{(k)})=\lim_kx^{(k+1)}=x^\ast$，所以$x^\ast$是不动点。这里三个要求各司其职：$T(X)\subseteq X$使迭代列有定义，统一的$q<1$使几何级数收敛，完备性使柯西列有极限。

\subsection{不动点的唯一性}
\label{sec:mmap-fixed-point-uniqueness}

设$x^\ast$与$y^\ast$都是不动点。由压缩条件，
\begin{equation*}
d(x^\ast,y^\ast)=d(T(x^\ast),T(y^\ast))\leq q\,d(x^\ast,y^\ast),
\end{equation*}
即$(1-q)\,d(x^\ast,y^\ast)\leq0$。由$q<1$得$d(x^\ast,y^\ast)=0$，两者相同。唯一性只用到$q<1$，不需要完备性：在不完备的空间中压缩映射可能没有不动点，但不会有两个。

结合上一小节，从任何初值出发的迭代都收敛到某个不动点，而不动点只有一个，所以所有迭代都趋于同一个$x^\ast$。也可以直接看出这一点：$d(x^{(k)},x^\ast)=d(T(x^{(k-1)}),T(x^\ast))\leq q\,d(x^{(k-1)},x^\ast)$，所以$d(x^{(k)},x^\ast)\leq q^k\,d(x^{(0)},x^\ast)$。误差每步至少缩小为原来的$q$倍，只是这个式子含有未知的$x^\ast$，不能直接计算。

\subsection{迭代误差界}
\label{sec:mmap-error-bounds}

式\eqref{eq:mmap-contraction-error-bounds}的两个界把未知的误差换成可以计算的量。

第一个界由式\eqref{eq:mmap-contraction-cauchy}令$m\to\infty$得到，因为$d(x^{(m)},x^{(k)})\to d(x^\ast,x^{(k)})$。它只用到初始步长$d(x^{(1)},x^{(0)})$，所以算出第一步后就能预先安排迭代次数：要使误差不超过$\tau$，只需
\begin{equation*}
\frac{q^k}{1-q}\,d(x^{(1)},x^{(0)})\leq\tau,\qquad\text{即}\qquad k\geq\frac{\ln\bigl(\tau(1-q)/d(x^{(1)},x^{(0)})\bigr)}{\ln q}.
\end{equation*}
这样的估计称为\term{先验误差界}（a priori error bound）。

第二个界利用已经算出的最近一步。由三角不等式与压缩条件，
\begin{equation*}
d(x^{(k)},x^\ast)\leq d(x^{(k)},x^{(k+1)})+d(x^{(k+1)},x^\ast)\leq q\,d(x^{(k)},x^{(k-1)})+q\,d(x^{(k)},x^\ast),
\end{equation*}
移项即得$(1-q)\,d(x^{(k)},x^\ast)\leq q\,d(x^{(k)},x^{(k-1)})$。同样的推导对任意点$x$给出$d(x,x^\ast)\leq d(x,T(x))/(1-q)$：一点的残差$d(x,T(x))$越小，它离不动点越近。这样的估计称为\term{后验误差界}（a posteriori error bound），可以作为停止判据：当$q\,d(x^{(k)},x^{(k-1)})/(1-q)\leq\tau$时停止，就保证误差不超过$\tau$。

两个界都依赖压缩常数，因子$1/(1-q)$在$q$接近$1$时很大。例如$q=0.99$时，后验界要求步长小于$\tau/99$才能保证误差小于$\tau$；同时误差每步只缩小$1\%$，收敛很慢。还要注意，后验判据的可靠性来自已知的$q$。不知道映射是否为压缩时，“步长已经很小”本身说明不了什么：调和级数的部分和每步只增加$1/k$，步长趋于零，部分和却趋于无穷。

\subsection{算例与反例}
\label{sec:mmap-contraction-examples}

先在一个可以完全手算的例子上核对定理。

\begin{example}[一次函数的迭代]
\label{ex:mmap-linear-iteration}
在完备空间$\mathbb R$上取$T(x)=(x+2)/3$。$\lvert T(x)-T(y)\rvert=\lvert x-y\rvert/3$，所以$T$是压缩常数$q=1/3$的压缩映射，唯一不动点由$x=(x+2)/3$解得$x^\ast=1$。从$x^{(0)}=0$出发，
\begin{equation*}
x^{(1)}=\tfrac23,\quad x^{(2)}=\tfrac89,\quad x^{(3)}=\tfrac{26}{27},\quad x^{(4)}=\tfrac{80}{81},
\end{equation*}
一般地$x^{(k)}=1-3^{-k}$，因为$1-x^{(k+1)}=(1-x^{(k)})/3$。实际误差是$3^{-k}$。

先验界为$\frac{(1/3)^k}{1-1/3}\cdot\frac23=3^{-k}$。取$k=3$，后验界为$\frac{1/3}{2/3}\bigl(\frac{26}{27}-\frac{24}{27}\bigr)=\frac1{27}$，也等于实际误差$3^{-3}$。两个界在这里都取等号，因为$T$是一次函数，压缩不等式本身就是等式。要使误差不超过$10^{-6}$，需要$3^{-k}\leq10^{-6}$，即$k\geq6\ln10/\ln3\approx12.6$，取$k=13$。
\end{example}

下面的例子说明定理的每个条件都不能省略。

缺少完备性时，极限可能不在空间中。在$X=(0,1)$上取$T(x)=x/2$，它把$X$映入自身，压缩常数为$1/2$。迭代$x^{(k)}=2^{-k}x^{(0)}$是柯西列，在$\mathbb R$中趋于$0$，但$0\notin X$；方程$x=x/2$在$\mathbb R$中只有解$0$，所以$T$在$X$中没有不动点。

压缩常数等于$1$时，距离不缩小，几何级数也不再收敛。在完备空间$\mathbb R$上取$T(x)=x+1$，有$\lvert T(x)-T(y)\rvert=\lvert x-y\rvert$，迭代每步移动$1$，不是柯西列；$x=x+1$无解。即使每对点的距离都严格缩小，只要没有统一的$q<1$，结论也可能失败。在完备空间$[1,\infty)$上取$T(x)=x+1/x$，它把$[1,\infty)$映入自身；由中值定理，$\lvert T(x)-T(y)\rvert=(1-1/\xi^2)\lvert x-y\rvert<\lvert x-y\rvert$（$x\neq y$，$\xi$介于两者之间），但$1-1/\xi^2$可以任意接近$1$。方程$x=x+1/x$无解，迭代列单调增加趋于无穷。

$T$必须把空间映回自身。在完备空间$[0,1/2]$上，$T(x)=(x+2)/3$仍满足压缩不等式，但$T(0)=2/3\notin[0,1/2]$，迭代离开了所给的空间，唯一的不动点$1$也不在其中。实际使用压缩映射定理时，常常先找一个被$T$映入自身的闭集，再在这个闭集上验证压缩条件；闭集作为完备空间的子空间仍完备。

压缩映射定理在其他地方多次出现。线性方程组的定常迭代$\bm x^{(k+1)}=\bm G\bm x^{(k)}+\bm c$中，若迭代矩阵在某个诱导范数下满足$\lVert\bm G\rVert<1$，迭代映射就是压缩常数为$\lVert\bm G\rVert$的压缩映射，这类格式的收敛分析可参考\BookChapterRef{chap:stationary-iterations}{卷二“定常迭代法”章}；常微分方程初值问题的Picard迭代在连续函数空间中验证压缩条件，由此得到解的存在唯一，可参考\BookChapterRef{chap:ordinary-differential-equations}{卷六“常微分方程与初值问题”章}。那里所用的完备函数空间，正是第~\ref{sec:mmap-uniform-cauchy}小节的有界连续函数空间。
